% Compile with: uplatex file.tex && dvipdfmx file.dvi
%           or: latexmk -uplatex file.tex
%\documentclass[uplatex,dvipdfmx]{jsarticle}
\documentclass[dvipdfmx,a4paper,11pt]{jsarticle}
\usepackage{amsmath, amssymb}
\usepackage{tikz}

\title{scenario\_ideal\_gas}
\date{\today}

\begin{document}
\maketitle

$N$ 粒子の単原子理想気体を扱う。粒子数 $N$、体積 $V$、温度 $T$、粒子の質量 $m$、プランク定数 $h$、ボルツマン定数 $k_B$ を正の実数パラメータとして準備する。

N, V, T, m, h, $k_B$ を正の実数パラメータとして定義して

\[
  N, V, T, m, h, k_B
\]

\noindent\textbf{Assumption:}\ $\displaystyle N > 0$\par

\noindent\textbf{Assumption:}\ $\displaystyle V > 0$\par

\noindent\textbf{Assumption:}\ $\displaystyle T > 0$\par

\noindent\textbf{Assumption:}\ $\displaystyle m > 0$\par

\noindent\textbf{Assumption:}\ $\displaystyle h > 0$\par

\noindent\textbf{Assumption:}\ $\displaystyle k_B > 0$\par

熱的ド・ブロイ波長 $\lambda=\sqrt{h^2/(2\pi mk_BT)}$（量子効果が効き始める長さスケール）を定義する。

$\lambda = \sqrt{h^2 / (2 \pi m k_B T)}$

\begin{equation}
  \lambda = \sqrt{h^2 / (2 \pi m k_B T)}
  \label{eq:E1}
\end{equation}

1粒子分配関数 $Z_1=V/\lambda^3$ を定義する（体積を熱的ド・ブロイ波長の3乗で割ったもので、粒子が占めうる量子状態の数に対応する）。

$Z_1 = V / E1^3$

\begin{equation}
  Z_1 = {{2^{\frac{3}{2}} \pi^{\frac{3}{2}} T^{\frac{3}{2}} V k_{B}^{\frac{3}{2}} m^{\frac{3}{2}}}\over{h^3}}
  \label{eq:E2}
\end{equation}

$N$ 粒子系の分配関数 $Z_N=Z_1^N/N!$ を定義する。$N!$ で割るのは、区別できない粒子の入れ替えによる数え過ぎを補正する Gibbs の補正項。

$Z_N = E2^N / N!$

\begin{equation}
  Z_N = {{2^{\frac{3 N}{2}} \pi^{\frac{3 N}{2}} T^{\frac{3 N}{2}} V^{N} k_{B}^{\frac{3 N}{2}} m^{\frac{3 N}{2}}}\over{N! h^{3 N}}}
  \label{eq:E3}
\end{equation}

ヘルムホルツ自由エネルギー $F=-k_BT\ln Z_N$ を定義する。

$F = -k_B T \log(E3)$

\begin{equation}
  F = -T k_{B} \log \left({{2^{\frac{3 N}{2}} \pi^{\frac{3 N}{2}} T^{\frac{3 N}{2}} V^{N} k_{B}^{\frac{3 N}{2}} m^{\frac{3 N}{2}}}\over{N! h^{3 N}}}\right)
  \label{eq:E4}
\end{equation}

これまでに登録した $Z_N,Z_1,\lambda$ の定義を代入し、$F$ を展開して $N,V,T,m,h,k_B$ だけで表される閉じた式にする。

F を展開して

\begin{equation}
  F = -T k_{B} \log \left({{2^{\frac{3 N}{2}} \pi^{\frac{3 N}{2}} T^{\frac{3 N}{2}} V^{N} k_{B}^{\frac{3 N}{2}} m^{\frac{3 N}{2}}}\over{N! h^{3 N}}}\right)
  \label{eq:E5}
\end{equation}

圧力 $P=-(\partial F/\partial V)_T$ を求めるため、まず $F$ を体積 $V$ で偏微分する。

F を V で偏微分して

\begin{equation}
  \frac{\partial F}{\partial V} = -\frac{N T k_{B}}{V}
  \label{eq:E6}
\end{equation}

符号を反転し、圧力 $P$ として登録する。

$P = -E6$

\begin{equation}
  P = \frac{N T k_{B}}{V}
  \label{eq:E7}
\end{equation}

得られた式を整理する。

E7 の右辺を整理して

\begin{equation}
  P = \frac{N T k_{B}}{V}
  \label{eq:E8}
\end{equation}

内部エネルギー $U=-T^2\partial(F/T)/\partial T$（Gibbs–Helmholtz の関係）を求めるため、まず $F/T$ を新しい式として作る。

$FoverT = F / T$

\begin{equation}
  FoverT = F / T
  \label{eq:E9}
\end{equation}

$F/T$ を温度 $T$ で偏微分する。

FoverT を T で偏微分して

\begin{equation}
  \frac{\partial FoverT}{\partial T} = -\frac{3 N k_{B}}{2 T}
  \label{eq:E10}
\end{equation}

$-T^2$ を掛けて内部エネルギー $U$ として登録する。

$U = -T^2 E10$

\begin{equation}
  U = \frac{3 N T k_{B}}{2}
  \label{eq:E11}
\end{equation}

得られた式を整理する。

E11 の右辺を整理して

\begin{equation}
  U = \frac{3 N T k_{B}}{2}
  \label{eq:E12}
\end{equation}

定積熱容量 $C_V=(\partial U/\partial T)_V$ を求めるため、内部エネルギー $U$ を温度 $T$ で偏微分する。

U を T で偏微分して

\begin{equation}
  \frac{\partial U}{\partial T} = \frac{3 N k_{B}}{2}
  \label{eq:E13}
\end{equation}

\end{document}